\writeSectionFrame{\presentationTitle}{Organisatorisches zur Vorlesung}
\frameWithImageOnly{\BILDERPFAD/organisatorisches1.jpg}
\note{
	Dann widmen wir uns erst einmal einigen organisatorischen Details.
}

\begin{frame}{Strukturierung der Vorlesung}
	\begin{itemize}
        \item Donnerstags, 1. Stunde (08:15 - 09:45 Uhr), Raum 4204
		\item Freitags, 1. Stunde (08:15 - 09:45 Uhr), Raum 3303
	\end{itemize}
\end{frame}

\begin{frame}{Themen}
	\begin{itemize}
		\item Was sind Entwurfsmuster?
        \item Kurze Einführung in UML
		\item Die traditionellen Entwurfsmuster der GoF
		\item Moderne Entwurfsmuster
	\end{itemize}
\end{frame}

\begin{frame}{Voraussetzungen}
	\begin{itemize}
		\item Muss: UML (vor allem Klassendiagramme)
		\item Kann: Java / C++ (da Beispiele in Java und C++)
	\end{itemize}
\end{frame}

\note{
	Es gibt zwei wichtige Voraussetzungen für dieses Modul. Sie müssen unbedingt Kenntnisse in UML haben, vor allem Klassendiagramm, da alle Muster sprachunabhängig mit UML-Diagrammen dargestellt werden. Sie müssen auch in der Prüfung in der Lage sein, ein Klassendiagramm zu einem Muster zu zeichnen und zu erklären. Kenntnisse in Java sind sinnvoll, weil die Beispiele alle in Java geschrieben sind. Hier ist es besonders wichtig, dass Sie sich in Vererbung, Polymorphie, Interfaces und abstrakten Klassen auskennen, denn auf diesen Mechanismen basieren beinahe alle Muster.  
}

\frameWithImageOnly{\BILDERPFAD/organisatorisches2.jpg}
\note{
 	Die Prüfung wird mündlich sein. Dabei sollten Sie in der Lage sein, verschiedene Entwurfsmuster zu erklären. Die Prüfung dauert ungefähr 20 Minuten. Kleiner Tipp: Bei Vorbereitung und Prüfung hilft es, sich eigene Beispiele zu den Mustern auszudenken. Dann hat man sie auch verstanden.
}

\frameWithImageOnly{\BILDERPFAD/organisatorisches3.jpg}
\note{
	Es gibt zu jedem Entwurfsmuster mehrere kleine Beispiele, die das Muster in verschiedenen Zusammenhängen zeigen. Es gibt allerdings auch noch ein großes Beispiel, in dem die Zusammenhänge zwischen den Mustern zu sehen sind und wie sie miteinander interagieren. Dieses Beispiel wird ein kleines textbasiertes Rollenspiel. 
}
